\documentclass[tikz,border=10pt]{standalone}
\usetikzlibrary{arrows.meta, positioning, shapes.geometric, shapes.misc, fit, backgrounds, calc, matrix}

\tikzset{
    >=Latex,
    font=\sffamily\large,
    node distance=1.5cm and 3cm,
    dag_result/.style={
        rectangle, 
        draw=green!70!black, 
        fill=green!15, 
        very thick, 
        minimum width=4.5cm, 
        minimum height=1.5cm, 
        text width=5cm,
        align=center
    },
    dag_lemma/.style={
        rounded rectangle, 
        draw=orange!80!black, 
        fill=yellow!20, 
        very thick, 
        minimum width=4.5cm, 
        minimum height=1.5cm, 
        text width=5cm,
        align=center
    },
    dag_model/.style={
        ellipse, 
        draw=blue!70!black, 
        fill=blue!12, 
        very thick, 
        minimum width=4.5cm, 
        minimum height=1.5cm, 
        text width=4.5cm,
        align=center
    },
    dag_unformalized/.style={
        rectangle, 
        draw=gray!60!black, 
        fill=gray!10, 
        very thick, 
        dashed, 
        minimum width=4.5cm, 
        minimum height=1.5cm, 
        text width=5cm,
        align=center
    },
    dag_conditional/.style={
        rectangle, 
        rounded corners,
        draw=orange!80!black, 
        fill=orange!10, 
        very thick, 
        minimum width=4.5cm, 
        minimum height=1.5cm, 
        text width=5cm,
        align=center
    },
    dag_partial/.style={
        rectangle,
        rounded corners,
        draw=yellow!55!black,
        fill=yellow!10,
        very thick,
        dashed,
        minimum width=4.5cm,
        minimum height=1.5cm,
        text width=5cm,
        align=center
    },
    dag_scaffold/.style={
        rectangle,
        draw=gray!55!black,
        fill=white,
        very thick,
        dotted,
        minimum width=4.5cm,
        minimum height=1.5cm,
        text width=5cm,
        align=center
    },
    dag_caveat/.style={
        diamond, 
        draw=red!70!black, 
        fill=red!10, 
        minimum width=4cm, 
        align=center, 
        inner sep=0pt, 
        font=\sffamily\small
    },
    dag_arrow/.style={
        -{Stealth[scale=1.2]}, 
        very thick,
        shorten >=1.5pt,
        shorten <=1.5pt
    },
    dag_dashed_arrow/.style={
        -{Stealth[scale=1.2]}, 
        very thick, 
        dashed,
        shorten >=1.5pt,
        shorten <=1.5pt
    },
    dag_template_legend/.style={
        minimum width=2.6cm,
        minimum height=0.8cm,
        text width=2.45cm,
        inner sep=3pt,
        font=\sffamily\small,
        align=center
    },
    dag_template_note/.style={
        rectangle,
        draw=gray!50,
        fill=gray!5,
        rounded corners,
        text width=13cm,
        align=left,
        font=\sffamily\small,
        inner sep=6pt
    }
}

% Legend Helper
\newcommand{\daglegend}[2]{
    \node[draw, rounded corners, very thick, fit=#1, inner sep=12pt, label=above:\textbf{\Large Legend}] (#2) {};
}


\begin{document}
\begin{tikzpicture}[
  x=1cm,
  y=1cm,
  every node/.append style={outer sep=5pt, font=\sffamily\small}
]

\dagPaperMetadata{Designing Optimal Binary Rating Systems}{Garg, Johari}{AISTATS / PMLR 89}{2019}{https://proceedings.mlr.press/v89/garg19a.html}

\dagPaperLegendRightOfMetadata{
\node[dag_model, dag_template_legend] (legDef) at (0,0) {definition /\\formula};
\node[dag_result, dag_template_legend] (legDone) at (4.2,0) {formalized\\result};
}
\daglegend{(legDef)(legDone)}{Legend}

\begin{dagPaperBody}

\node[dag_model] (KL) at (0,0) {
  \textbf{Bernoulli KL (Sec. 3.3)} \\
  finite-support KL and support-safe convention
};

\node[dag_result] (AdjRate) at (6.2,0) {
  \textbf{Theorem 3.1 Rate Formula} \\
  adjacent binary threshold rate in real and support-safe forms
};

\node[dag_result] (ClosedRate) at (12.4,0) {
  \textbf{Lemma 3.1 Closed Rate} \\
  weighted common-threshold expression realizes the displayed value
};

\node[dag_result] (LevelOpt) at (18.6,0) {
  \textbf{Lemma 3.1 Finite Levels} \\
  endpoint-aware equalized levels exist, are unique, and are maximin optimal
};

\node[dag_result] (T31Lex) at (24.8,0) {
  \textbf{Theorem 3.1 Lexicographic Optimality} \\
  value first, then rate, characterizes optimal binary systems
};

\node[dag_model] (MatchModel) at (0,-3.3) {
  \textbf{Monotone Match Model} \\
  nondecreasing match-function primitive used by appendix rates
};

\node[dag_result] (C1C2) at (6.2,-3.3) {
  \textbf{Lemmas C.1-C.2} \\
  pairwise binary error probabilities have the Bernoulli exponent
};

\node[dag_result] (FiniteRate) at (12.4,-3.3) {
  \textbf{Finite Adjacent Objective} \\
  endpoint/source-shaped rows construct finite certificates and give exact rates
};

\node[dag_result] (T31Branches) at (18.6,-3.3) {
  \textbf{Theorem 3.1 Fixed/Continuum Branches} \\
  fixed discretizations and continuum cell integrals realize the same optimum
};

\node[dag_result] (T31) at (24.8,-3.3) {
  \textbf{Full Theorem 3.1} \\
  continuum optimal discretization and maximin decomposition are checked
};

\node[dag_result] (TC1) at (0,-6.6) {
  \textbf{Theorem C.1} \\
  weighted Laplace principle converts local rates to integral rates
};

\node[dag_result] (C3) at (6.2,-6.6) {
  \textbf{Lemma C.3} \\
  piecewise-constant partitions aggregate by the minimum adjacent exponent
};

\node[dag_result] (C4) at (12.4,-6.6) {
  \textbf{Lemma C.4} \\
  positive rate occurs exactly for the source piecewise-constant case
};

\node[dag_result] (C5) at (0,-9.9) {
  \textbf{Lemma C.5} \\
  doubled endpoint chain is feasible and equalizes adjacent rates
};

\node[dag_result] (C1C2C6C8) at (6.2,-9.9) {
  \textbf{Cor. C.1-C.2; Lemmas C.6-C.8} \\
  rate vanishes, mesh shrinks, and finite-support bounds hold
};

\node[dag_result] (C3Cor) at (12.4,-9.9) {
  \textbf{Corollary C.3} \\
  monotone first-level lower bound supports the finite grid guarantee
};

\node[dag_result] (C9) at (18.6,-9.9) {
  \textbf{Lemma C.9} \\
  nested-bisection operation count satisfies the step bound
};

\node[dag_result] (T32Cert) at (24.8,-9.9) {
  \textbf{Theorem 3.2} \\
  calculated grid gives additive loss and runtime bounds
};

\node[dag_result] (B1Lemma) at (0,-13.2) {
  \textbf{Lemma B.1} \\
  matching-rate shifts preserve the limiting comparison
};

\node[dag_result] (TB1) at (6.2,-13.2) {
  \textbf{Theorem B.1} \\
  quantile-floor representation and uniform convergence give uniform limits
};

\node[dag_result] (B2B3) at (12.4,-13.2) {
  \textbf{Lemmas B.2-B.3} \\
  known-type and unknown-type learning converge uniformly
};

\node[dag_model] (DefC1) at (18.6,-13.2) {
  \textbf{Definition C.1} \\
  Kendall and Spearman population objectives
};

\node[dag_result] (C10C12) at (24.8,-13.2) {
  \textbf{Lemmas C.10-C.12} \\
  Spearman reduction and Kendall/Spearman objectives favor equispaced cutpoints
};

\node[dag_result] (CorC4) at (24.8,-16.5) {
  \textbf{Corollary C.4} \\
  equispaced Kendall/Spearman rules arise along a uniform optimal subsequence
};

\draw[dag_arrow] (KL.east) -- (AdjRate.west);
\draw[dag_arrow] (AdjRate.east) -- (ClosedRate.west);
\draw[dag_arrow] (ClosedRate.east) -- (LevelOpt.west);
\draw[dag_arrow] (LevelOpt.east) -- (T31Lex.west);

\draw[dag_arrow] (KL.south east) -- (C1C2.north west);
\draw[dag_arrow] (AdjRate.south) -- (C1C2.north);
\draw[dag_arrow] (C1C2.east) -- (FiniteRate.west);
\draw[dag_arrow] (FiniteRate.east) -- (T31Branches.west);
\draw[dag_arrow] (T31Branches.east) -- (T31.west);
\draw[dag_arrow] (T31Lex.south) -- (T31.north);

\draw[dag_arrow] (MatchModel.south east) -- (C3.north west);
\draw[dag_arrow] (TC1.east) -- (C3.west);
\draw[dag_arrow] (C3.east) -- (C4.west);
\draw[dag_arrow] (C4.north east) -- (T31Branches.south west);

\draw[dag_arrow] (C5.east) -- (C1C2C6C8.west);
\draw[dag_arrow] (C1C2C6C8.east) -- (C3Cor.west);
\draw[dag_arrow] (C3Cor.east) -- (C9.west);
\draw[dag_arrow] (C9.east) -- (T32Cert.west);

\draw[dag_arrow] (B1Lemma.east) -- (TB1.west);
\draw[dag_arrow] (C1C2C6C8.south) to[out=-90,in=115,looseness=1.05] (TB1.north west);
\draw[dag_arrow] (TB1.east) -- (B2B3.west);

\draw[dag_arrow] (DefC1.east) -- (C10C12.west);
\draw[dag_arrow] (C10C12.south) -- (CorC4.north);
\draw[dag_arrow] (TB1.south east) to[out=-30,in=180,looseness=1.1] (CorC4.west);

\end{dagPaperBody}
\end{tikzpicture}
\end{document}
